\section{CMC sphere jets and their Hawking mass}
\label{ads:sec:cmc}

We use the finite Taylor conventions of Section~\ref{ads:sec:taylor}.
Taylor expansions of inverses, compositions, and other smooth geometric
operations at a rotation-invariant background preserve finite angular
type at each fixed order: each coefficient uses only finitely many
linear differential operations and tensor products.

\subsection{The sphere construction}

In the end put \(\eta=\operatorname{arsinh}r\), so that
\[
 b=\dd\eta^2+\sinh^2\eta\,\sigma.
\]
An estimate for a tensor below refers to its components in scaled
coframes \(\dd\eta,r\,\dd x^A\), in fixed smooth angular coordinate
charts.  A differentiated little-oh estimate includes all mixed
radial and angular derivatives of the stated total order.

\begin{proposition}[CMC sphere jets]\label{ads:prop:cmc-jets}
Let \(p\geq1\), and let \(G\) be a smooth Riemannian metric jet of order
\(p\), with base \(G_0=g_0\), whose coefficients have finite angular type.
Suppose that there are angular scalar jets \(C,E\), with
\(C_0=2m_0\) and \(E_0=0\), such that
\begin{equation}\label{ads:cmc:metric-asymptotics}
 G-b=r^{-3}\bigl(C\,\dd\eta^2+E r^2\sigma\bigr)+o(r^{-3}).
\end{equation}
For the coefficient of degree \(j\), \(1\leq j\leq p\), suppose that
the remainder satisfies this estimate after every mixed radial and
angular derivative of total order at most \(p-j+1\).

There is a smooth family of sphere jets
\(\Sigma_s\), \(0\leq s<\infty\), given by graphs over
\(\{s\}\times S^2\), with the following properties:
\begin{enumerate}
 \item \(\Sigma_s\) has constant outward mean curvature through order
 \(p\), and \(\Sigma_0\) is minimal through that order.
 \item The graph coefficients have finite angular type.  The normal
 lapse \(f\) for variation in \(s\) has base value \(f_0=1\).
 \item If \(A(s)\) and \(H(s)\) are the area and mean curvature jets,
 then the CMC Hawking mass
 \begin{equation}\label{ads:cmc:hawking-definition}
 m_H(\Sigma_s)=
 \sqrt{\frac{A(s)}{16\pi}}
 \left(1-\frac{(H(s)^2-4)A(s)}{16\pi}\right)
 \end{equation}
 satisfies, coefficientwise,
 \begin{equation}\label{ads:cmc:hawking-limit}
 \lim_{s\to\infty}m_H(\Sigma_s)
 =m_{\AH}(G)
 :=\sqrt{p_0(G)^2-\sum_{i=1}^3p_i(G)^2}.
 \end{equation}
 Here \(p_\mu(G)\) is the coefficientwise flux \eqref{ads:eq:mass-flux},
 and the square root is its formal branch with base value
 \(m_0>0\).
\end{enumerate}
\end{proposition}

The expression \eqref{ads:cmc:hawking-definition} is the CMC specialization
of the asymptotically hyperbolic Hawking mass used in
\cite[Equation~(2)]{Ambrozio2015}.  The proof below checks its limiting
normalization against \eqref{ads:eq:mass-flux} directly.

\begin{proof}
We first solve the Jacobi equations on finite radial intervals.
At infinity, hyperbolic boosts account for the order-one degree-one
displacement. The asymptotic CMC equation selects a frame in which the
spatial mass components vanish. We then choose the free graph means
to join those spheres to the minimal first leaf.

\paragraph{The graph equation on finite radial intervals.}
For a normal variation with speed \(w\), the mean-curvature variation is
\begin{equation}\label{ads:cmc:jacobi-general}
 Jw=-\Delta_{\Sigma}w
 -\bigl(\Ric_G(\nu,\nu)+|\II|^2\bigr)w.
\end{equation}
Indeed, the induced metric varies by \(2w\II\), the variation of the
unit normal is \(-\grad_\Sigma w\), and differentiation of the second
form gives the Hessian term, the ambient curvature term, and the
quadratic second-form term.  Tracing, including the variation of the
inverse induced metric, gives \eqref{ads:cmc:jacobi-general}.  A tangential
velocity adds the directional derivative of \(H\).

For a background coordinate sphere, \(H_0=2N/r\) and
\[
 \Ric_{g_0}(\partial_s,\partial_s)=-2-\frac{2m_0}{r^3},
 \qquad |\II_0|^2=\frac{2N^2}{r^2}.
\]
Thus its Jacobi operator is
\begin{equation}\label{ads:cmc:background-jacobi}
 J_s^0=-\Delta_{r^2\sigma}-\frac2{r^2}+\frac{6m_0}{r^3}.
\end{equation}
On degree-\(\ell\) spherical harmonics its eigenvalue is
\begin{equation}\label{ads:cmc:jacobi-eigenvalues}
 \lambda_\ell(s)
 =\frac{\ell(\ell+1)-2}{r^2}+\frac{6m_0}{r^3}.
\end{equation}
It is strictly positive on every nonconstant mode.  In particular,
the degree-one eigenvalue is \(6m_0/r^3>0\).
At \(s=0\), using \(2m_0=a(1+a^2)\), the constant eigenvalue is
\[
 -\frac2{a^2}+\frac{6m_0}{a^3}
 =3+\frac1{a^2}=\frac{2\kappa}{a},
\]
and hence \(J_0^0=D\), positive on all modes.

Write a graph as \(s+\rho(s,x;\eps)\), with \(\rho_0=0\).
At degree \(j\), its mean-curvature coefficient is
\[
 J_s^0\rho_j+\mathcal R_j(s,x),
\]
where \(\mathcal R_j\) is known from the metric coefficients and the
lower graph coefficients.  To make the first sphere minimal, solve
this equation at \(s=0\) with right-hand side zero, using \(D^{-1}\).
At general \(s\), prescribe any smooth round mean for \(\rho_j\) that
agrees with this value at zero, and solve the projection of the
equation onto the nonconstant modes.
The operator \eqref{ads:cmc:background-jacobi} preserves the splitting
into constants and nonconstants, so the latter equation has a unique
solution.  At each order it is a finite-dimensional linear equation
with a smooth invertible coefficient matrix at every finite radius.
This produces smooth coefficients on each finite interval.
Induction also shows that the resulting graph at zero is exactly the
minimal graph jet already constructed there.
The constant eigenvalue need not be inverted away from zero.

\paragraph{The flux aspect.}
For an angular scalar \(F\), write
\(\langle F\rangle_\sigma=(4\pi)^{-1}\int_{S^2}F\,\dd A_\sigma\).
The asymptotic form \eqref{ads:cmc:metric-asymptotics} gives
\begin{equation}\label{ads:cmc:mass-aspect}
 (p_0(G),p_1(G),p_2(G),p_3(G))
 =\left\langle\mu(x)(1,x)\right\rangle_\sigma,
 \qquad \mu=\frac{C+3E}{2}.
\end{equation}
Here \(\mu\) denotes only a flux aspect.
To verify the formula directly, let \(e=G-b\).  Its leading trace is
\((C+2E)r^{-3}\), and its radial divergence-minus-trace derivative is
\begin{align*}
 (\Div_b e-\dd\tr_b e)(\partial_\eta)
 &=2\coth\eta(C-E)r^{-3}
       -2\partial_\eta(Er^{-3})+o(r^{-3})\\
 &=(2C+4E)r^{-3}+o(r^{-3}).
\end{align*}
For \(V_0=\sqrt{1+r^2}\) or \(V_i=rx_i\), let \(v=1\) or \(x_i\),
respectively.  Both \(V\) and \(\partial_\eta V\) have leading term
\(rv\).  The terms \((\tr_b e)\dd V-e(\grad_bV,\cdot)\) add
\(2Ev r^{-2}+o(r^{-2})\) to the normal flux integrand.
The total is therefore
\((2C+6E)v r^{-2}+o(r^{-2})\).
Multiplying by \(r^2\dd A_\sigma/(16\pi)\) proves
\eqref{ads:cmc:mass-aspect}.  The first differentiated remainder assumed
in \eqref{ads:cmc:metric-asymptotics} makes every omitted flux integral
tend to zero.

\paragraph{Boosted coordinates.}
Realize hyperbolic space as
\[
 (\cosh\eta,\sinh\eta\,x)\subset\mathbb R^{3,1}.
\]
Let \(\mathcal B(d)\) be the exponential of the Lorentz Lie-algebra
matrix with time-space entries \(d\in\mathbb R^3\) and zero spatial
rotation part.  Thus \(\mathcal B(0)\) is the identity.
Write its action on future null vectors as
\begin{equation}\label{ads:cmc:boost-null}
 \mathcal B(d)(1,x)=k(x)(1,y(x)).
\end{equation}
For small \(d\), \(k>0\).
In pullback coordinates the old radius is \(kr+O(r^{-1})\), the old
radial distance is \(\eta+\log k+O(r^{-2})\), and the old radial unit
covector agrees with the new one up to \(O(r^{-1})\) in hyperbolic
norm.  These statements, including all their needed derivatives,
follow by applying the matrix \(\mathcal B(d)\) to the displayed
hyperboloid coordinates.  Since \(\mathcal B(d)\) preserves \(b\),
the leading coefficients of its pullback of \(G\) are
\begin{equation}\label{ads:cmc:boost-aspect}
 C_d=k^{-3}C\circ y,\qquad E_d=k^{-3}E\circ y.
\end{equation}
For example, the angular part of \(\dd\log k\) has hyperbolic norm
\(O(r^{-1})\); it therefore changes neither leading scaled diagonal
coefficient in \eqref{ads:cmc:metric-asymptotics}.

Differentiating \eqref{ads:cmc:boost-null} tangentially and taking
Minkowski inner products gives \(y^*\sigma=k^{-2}\sigma\).
Consequently \(\dd A_\sigma(y)=k^{-2}\dd A_\sigma(x)\), and
\((1,x)=k\mathcal B(d)^{-1}(1,y)\).  Changing angular variables proves
the exact leading-flux transformation
\begin{equation}\label{ads:cmc:boost-mass}
 \left\langle k^{-3}\mu(y)(1,x)\right\rangle_\sigma
 =\mathcal B(d)^{-1}\left\langle\mu(x)(1,x)\right\rangle_\sigma.
\end{equation}
This establishes the needed covariance directly for the specified
flux.

\paragraph{The asymptotic CMC equation.}
For large \(s\), put \(R=r(s)\) and \(\eta_R=\operatorname{arsinh}R\).
In the coordinates obtained from \(\mathcal B(d)\), seek a graph
\begin{equation}\label{ads:cmc:far-graph}
 \eta=\eta_R+\frac{h(x)}{R},
\end{equation}
where \(d\) and \(h\) are jets with zero base values, and every
coefficient of \(h\) has zero constant and degree-one components.
The parameters \(d,h\) are held fixed when computing an individual
leaf.  We claim, coefficientwise for bounded input jets in fixed
finite angular spaces, that
\begin{equation}\label{ads:cmc:mean-curvature-asymptotic}
 R^3\bigl(H-2\coth\eta_R\bigr)
 =-(\Delta_\sigma+2)h-(C_d+3E_d)+o(1).
\end{equation}

Here are the estimates underlying this claim.
In \(b\), the unit normal to \eqref{ads:cmc:far-graph} is
\[
 \frac{\partial_\eta-r^{-2}\grad_\sigma(h/R)}
 {\sqrt{1+r^{-2}|\dd(h/R)|_\sigma^2}}.
\]
Its divergence gives
\[
 H_b=2\coth\eta_R
       -R^{-3}(\Delta_\sigma+2)h+O(R^{-4}).
\]
Indeed the slope has hyperbolic norm \(O(R^{-2})\), its normalization
differs from one by \(O(R^{-4})\), and the second Taylor term in
\(2\coth(\eta_R+h/R)\) is \(O(R^{-4})\).
These estimates hold for every coefficient of the finite Taylor
expansion when the input coefficients are bounded.

On an unshifted sphere the perturbation in
\eqref{ads:cmc:metric-asymptotics}, after boosting, changes the radial
lapse by \(C_dR^{-3}/2+o(R^{-3})\).  This contributes
\(-C_dR^{-3}+o(R^{-3})\) to \(H\).
The relative area density changes by \(E_dR^{-3}+o(R^{-3})\);
its radial logarithmic derivative contributes
\(-3E_dR^{-3}+o(R^{-3})\).
The mixed components in the remainder contribute \(o(R^{-3})\),
including their tangential-divergence term.
For completeness, on a fixed surface the comparison between two
metrics \(b\) and \(\widehat g\) is
\[
 \II_{\widehat g}(X,Y)
 =\widehat g(\nu_{\widehat g},\nu_b)\II_b(X,Y)
 -\widehat g\bigl(\nu_{\widehat g},
       (\nabla^{\widehat g}-\nabla^b)_XY\bigr).
\]
It follows by splitting \(\nabla_X^bY\) into its tangential and normal
parts and using normal orthogonality.
After tracing, the linear metric correction depends smoothly on
the tangent plane, the background shape operator, and the scaled
components of \(\widehat g-b,\nabla^b(\widehat g-b)\).
The slope of \eqref{ads:cmc:far-graph} tends to zero, and its background
shape operator differs by \(o(1)\) from that of a coordinate sphere.
Its displacement is \(O(R^{-1})\) in radial distance.
The leading metric correction on that graph is consequently still
\(-(C_d+3E_d)R^{-3}\).
Quadratic metric corrections are \(O(R^{-6})\).
This proves \eqref{ads:cmc:mean-curvature-asymptotic}.

The estimate just proved also justifies its use at every Taylor
order through \(p\).
A coefficient of \(G\) of degree \(j\) can undergo at most \(p-j\)
spatial differentiations from Taylor composition with the graph and
the boost.  Mean curvature requires one additional derivative of
the metric.  Thus the mixed derivative allowance \(p-j+1\) in
\eqref{ads:cmc:metric-asymptotics} controls the coefficient of every
remainder after multiplication by \(R^3\).
The base \(g_0\) has the required bounds to all orders.
Angular differential operators cause no further restriction within
a fixed finite-dimensional angular space.
More explicitly, denote the degree-\(j\) residual in
\eqref{ads:cmc:metric-asymptotics} by \(\mathcal E_j\), and set
\[
 \omega_j(R)=
 \max_{a+b\le p-j+1}\ 
 \sup_{r\ge R/2}
 \left\|\partial_\eta^a\nabla_\sigma^b
              (r^3\mathcal E_j^{\mathrm{scaled}})\right\|_\infty.
\]
Then \(\omega_j(R)\to0\).
For every \(n\le p\), bounded coefficients of \(d,h\) in fixed finite
angular spaces give the estimate
\[
 \left\|[\,
 R^3(H-2\coth\eta_R)+(\Delta_\sigma+2)h+C_d+3E_d\,]_n
       \right\|_\infty
 \le C_{p,\mathcal V,M}
       \left(R^{-1}+\sum_{j=1}^n\omega_j(R)\right).
\]
Here \(\mathcal V\) denotes the finitely many angular spaces and \(M\)
a bound for the input coefficients; the fixed metric coefficients
may also enter the constant.  The explicit base residual is
\(O(r^{-5})\) with all derivatives and is included in the
\(R^{-1}\) bound.  The preceding graph estimates, the product rule,
and the derivative count prove this estimate term by term.
It is an estimate for Taylor coefficients, exactly as needed for
the recursion.

\paragraph{Solving the rescaled equation.}
Project \(R^3(H-2\coth\eta_R)\) onto the nonconstant modes and set the
projection equal to zero, through order \(p\).
At the background, the finite-radius linearization preserves the
splitting between the \(h\)-modes, of degrees at least two, and the
three boost modes.
Set
\[
 \alpha_R=\frac{\sqrt{1+R^2}}{N(R)}.
\]
A graph variation \(h/R\) in \(\eta\) has physical normal
displacement \(\alpha_Rh/R\).  Equations
\eqref{ads:cmc:background-jacobi}--\eqref{ads:cmc:jacobi-eigenvalues} give its
exact degree-\(\ell\) multiplier in the rescaled equation:
\begin{equation}\label{ads:cmc:scaled-inverse-blocks}
 \alpha_R\bigl(\ell(\ell+1)-2+6m_0/R\bigr),
 \qquad \ell\ge2.
\end{equation}
An infinitesimal boost has \(\eta\)-displacement \(d\cdot x\).
Its physical normal displacement is \(\alpha_Rd\cdot x\), so its
multiplier is \(6m_0\alpha_R\).
These blocks converge respectively to \(\ell(\ell+1)-2\) and
\(6m_0\), both invertible.
Equivalently, the latter sign and coefficient follow by
differentiating \(2m_0k^{-3}\) in
\eqref{ads:cmc:mean-curvature-asymptotic}, since
\(k=1+d\cdot x+O(|d|^2)\).

At each Taylor order only fixed finite angular spaces are involved.
In particular, the Taylor coefficients at \(d=0\) of the boost map
and its compositions are finite sums of coordinate functions and
angular derivatives of the given profiles.  This assertion concerns
the finite Taylor expansion: it does not require a finite harmonic
cutoff to be preserved by a finite nonzero boost.
One can see the closure directly from the infinitesimal boost field
\[
 \mathcal K_a=(a\cdot x)\partial_\eta
               +\coth\eta\,\grad_\sigma(a\cdot x).
\]
It follows by differentiating the hyperboloid coordinates.
Its Lie derivative, repeated finitely many times, uses only radial
derivatives, angular derivatives, and multiplication by degree-one
profiles, also for tensor fields.
It follows that the inverses of the finite-dimensional blocks above
are bounded for all sufficiently large \(R\) and converge as
\(R\to\infty\).

Now solve successively for the coefficients of \(h,d\).
At each order their linear operator is the displayed background
linearization.  The remaining source depends only on lower solved
coefficients and the prescribed metric coefficients.
Equation \eqref{ads:cmc:mean-curvature-asymptotic}, with its uniform
coefficientwise remainder estimate, shows inductively that this
source converges once the lower coefficients do.
The converging inverse then gives bounded coefficients of \(h,d\)
with limits.  This starts at degree one and proves the assertion
through order \(p\).
Smooth dependence on \(s\) follows at every finite radius from the
smooth finite-dimensional equations.

In original coordinates these spheres are graph jets over the
background sphere: the angular map has identity as base and can be
inverted formally.  Their round mean graph heights are smooth scalar
functions of \(s\).  Choose the free means in the finite-radius
recursion to agree with these functions for all sufficiently large
\(s\), and to agree with the minimal graph's means at zero, using a
smooth interpolation.  Inductive uniqueness of the nonconstant
graph coefficients makes the resulting recursion coincide with the
far construction.  This gives the family on the whole half-line.
Its base is the coordinate-sphere family, so its lapse has base one.

\paragraph{The Hawking mass limit.}
Let \(d_\infty\) be the coefficientwise limiting boost.
The degree-one projection of
\eqref{ads:cmc:mean-curvature-asymptotic} gives
\[
 \left\langle\mu_{d_\infty}x_i\right\rangle_\sigma=0,
 \qquad
 \mu_d=(C_d+3E_d)/2.
\]
Writing
\[
 H_b(R)=2\coth\eta_R,\qquad
 B_R=R^3(H-H_b(R)),
\]
the round average of that same equation gives
\[
 B_R=-2\langle\mu_d\rangle_\sigma+o(1).
\]
The area of \eqref{ads:cmc:far-graph} satisfies
\begin{equation}\label{ads:cmc:area-error}
 \frac{A}{4\pi R^2}=1+\delta_R,\qquad \delta_R=O(R^{-2})
\end{equation}
coefficientwise.  To see the relevant cancellation, the relative
hyperbolic area density is
\[
 1+2\coth\eta_R\,\frac hR+O(R^{-2}),
\]
and the metric correction is \(O(R^{-3})\).  The sole possible
relative term of order \(R^{-1}\) integrates to zero because every
coefficient of \(h\) has zero round mean.

Since \(H_b(R)^2-4=4/R^2\), direct substitution into
\eqref{ads:cmc:hawking-definition} gives
\begin{align*}
 m_H={}&-\frac{R\delta_R}{2}\sqrt{1+\delta_R}
 -\frac{H_b(R)B_R}{4}(1+\delta_R)^{3/2}\\
 &-\frac{B_R^2}{8R^3}(1+\delta_R)^{3/2}.
\end{align*}
Therefore
\[
 \lim_{s\to\infty}m_H(\Sigma_s)
 =\langle\mu_{d_\infty}\rangle_\sigma.
\]
By \eqref{ads:cmc:boost-mass}, this last scalar is the time component of
\(\mathcal B(d_\infty)^{-1}p(G)\), whose spatial components vanish.
Lorentz invariance and its positive base value \(m_0\) identify it
with the formal square root in \eqref{ads:cmc:hawking-limit}.
\end{proof}

\begin{remark}\label{ads:cmc:regularity-budget}
Lemma~\ref{ads:ana:coefficients} supplies the derivative budget used here:
for a general finite-angular-type seed, the homogeneous first
coefficient has all differentiated asymptotics and the second has
the first differentiated asymptotic required at order two.
The fourth-order construction is used only after reduction to
$q=\Hess v-Bv$, with $Lv=0$ and $v$ consisting of constant and
degree-one angular modes. All required differentiated decay is
available in this class.
\end{remark}

\subsection{The Hawking variation identity}

On the sphere jets of Proposition~\ref{ads:prop:cmc-jets}, write
\[
 \langle u\rangle=\frac1A\int_{\Sigma_s}u\,\dd A_G,
 \qquad \bar f=\langle f\rangle,\qquad \delta f=f-\bar f.
\]
Here averages are with respect to the induced metric; the round
average retains the subscript \(\sigma\).
Let \(J\) be the full Jacobi operator
\eqref{ads:cmc:jacobi-general}, and let \(\II^\circ\) be the tracefree
second fundamental form.

\begin{lemma}[CMC Hawking variation]\label{ads:lem:hawking-variation}
The following identity holds through the order of the sphere and
metric jets:
\begin{equation}\label{ads:cmc:hawking-variation}
 \frac{\dd}{\dd s}m_H
 =\sqrt{\frac A{16\pi}}\frac{HA}{8\pi}
 \left[
 \bar f\left\langle
      \frac{\scal_G+6+|\II^\circ|^2}{2}
       \right\rangle
 +\frac1{\bar f}\langle\delta f\,J\delta f\rangle
 \right].
\end{equation}
\end{lemma}

\begin{proof}
The area and mean-curvature variation formulas give
\[
 A'=AH\bar f,\qquad H'=Jf.
\]
The latter is a constant on each leaf; any tangential transport
term vanishes since \(H\) is constant there.
For brevity put
\[
 c=\frac{4\pi}{A}+3-\frac{3H^2}{4},
 \qquad
 Q_H=\frac{\scal_G+6+|\II^\circ|^2}{2}.
\]
Differentiation of \eqref{ads:cmc:hawking-definition} gives
\begin{equation}\label{ads:cmc:raw-variation}
 m_H'=\sqrt{\frac A{16\pi}}\frac{HA}{8\pi}
          (c\bar f-H').
\end{equation}
The Gauss equation and
\(|\II|^2=H^2/2+|\II^\circ|^2\) imply
\[
 J1=K_\Sigma+3-\frac{3H^2}{4}-Q_H,
\]
where \(K_\Sigma\) is intrinsic Gauss curvature.
The sphere jets have \(\int K_\Sigma\,\dd A_G=4\pi\)
coefficientwise.
One can verify this without a limiting argument: the variation of
integrated scalar curvature on a closed two-dimensional metric is
the integral of a divergence minus its variation tensor paired with
\(\Ric-\frac12\scal\,g\).  The latter tensor is zero in dimension
two, so that integral is constant under every variation and equals
its round value.  Applying the same identity to the Taylor
coefficients proves the assertion for jets.
Consequently
\[
 \langle J1\rangle=c-\langle Q_H\rangle.
\]

The operator \(J\) is self-adjoint on the closed leaf.
Since \(\langle\delta f\rangle=0\) and \(Jf=H'\) is constant,
\begin{align*}
 H'&=\bar f\bigl(c-\langle Q_H\rangle\bigr)
                  +\langle\delta f\,J1\rangle,\\
 0&=\langle\delta f\,Jf\rangle
   =\bar f\langle\delta f\,J1\rangle
                  +\langle\delta f\,J\delta f\rangle.
\end{align*}
The scalar jet \(\bar f\) is invertible because its base value is one.
Eliminating the common term gives
\[
 c\bar f-H'
 =\bar f\langle Q_H\rangle
       +\bar f^{-1}\langle\delta f\,J\delta f\rangle.
\]
Substitution into \eqref{ads:cmc:raw-variation} proves the identity.
All calculations may first be made on a compact radial interval.
Smooth representatives of the metric and graph jets realize the
ordinary geometric variation identities there, and the imposed CMC
identities hold to the required Taylor order.
\end{proof}

\subsection{The first nonzero coefficient}

\begin{proposition}[Leading coefficient as a positive integral]
\label{ads:prop:leading-energy}
Use the construction of Proposition~\ref{ads:prop:cmc-jets} through order
\(2l\), where \(l\ge1\), and let \(A_{\mathrm{leaf}}=A(0)\) be the area
jet of its minimal first leaf.
Suppose, as ambient and leafwise coefficient identities, respectively,
that
\begin{equation}\label{ads:cmc:leading-assumptions}
 \scal_G+6=\eps^{2l}|k|_{g_0}^2+O(\eps^{2l+1}),
 \qquad
 \II^\circ=\eps^lU+O(\eps^{l+1}),
 \qquad
 \delta f=\eps^lF+O(\eps^{l+1}).
\end{equation}
Here \(k\) is a smooth tensor on \(M\); the fields \(U,F\) on the
background spheres are identified with fields on \(M\).
Then every coefficient below degree \(2l\) of
\(m_{\AH}(G)-b_*(A_{\mathrm{leaf}})\) vanishes, and
\begin{equation}\label{ads:cmc:leading-integral}
 \begin{split}
 &\bigl[m_{\AH}(G)-b_*(A_{\mathrm{leaf}})\bigr]_{2l}\\
 &\quad=\frac1{16\pi}\int_0^\infty N(s)
   \int_{\{s\}\times S^2}
    \left(|k|^2+|U|^2+2FJ_s^0F\right)\dd A_{g_0}\,\dd s
 \ge0.
 \end{split}
\end{equation}
All norms in this formula use the background metrics.
The displayed iterated integral is finite.
If it vanishes, then \(k,U,F\) vanish for \(s>0\), and also at the
boundary by continuity.
\end{proposition}

\begin{proof}
The first nonzero coefficient \(F\) has zero round mean on each
background sphere: extract degree \(l\) from
\(\int_{\Sigma_s}\delta f\,\dd A_G=0\), using the vanishing of all
lower coefficients of \(\delta f\).
Every term in the bracket in \eqref{ads:cmc:hawking-variation} starts at
degree \(2l\).  At that degree, all exterior factors can therefore
be evaluated at the background:
\[
 A_0=4\pi r^2,\quad H_0=2N/r,\quad \bar f_0=1,
 \quad\sqrt{\frac{A_0}{16\pi}}H_0=N.
\]
The same reasoning evaluates the norms, the induced measure, and
the Jacobi operator inside the leading terms at the background.
Composing the ambient scalar coefficient with a graph displacement
changes only higher degrees.
It follows that
\begin{equation}\label{ads:cmc:leading-derivative}
 [m_H']_{2l}
 =\frac{N(s)}{16\pi}\int_{\{s\}\times S^2}
        \bigl(|k|^2+|U|^2+2FJ_s^0F\bigr)\dd A_{g_0},
\end{equation}
and all lower coefficient derivatives vanish.

The integrand in \eqref{ads:cmc:leading-derivative} is nonnegative after
sphere integration.  Indeed, on the mean-zero field \(F\),
\[
 \int_{\{s\}\times S^2}FJ_s^0F\,\dd A_{g_0}
 =\sum_{\ell\ge1}
  \left(\frac{\ell(\ell+1)-2}{r^2}+\frac{6m_0}{r^3}\right)
       \|F_\ell\|_{L^2(r^2\sigma)}^2.
\]
Only finitely many modes occur, and every displayed eigenvalue is
strictly positive at a finite radius.
Also \(N(s)>0\) for \(s>0\).

For any finite \(T\), integrate the coefficient identities on
\([0,T]\).
The first leaf is minimal to the required order, so
\(m_H(\Sigma_0)=b_*(A_{\mathrm{leaf}})\) through that order.
The coefficientwise endpoint limits exist by
\eqref{ads:cmc:hawking-limit}.
For degree \(2l\), the right-hand side of the integrated
\eqref{ads:cmc:leading-derivative} is an increasing function of \(T\)
and its limit equals the finite limiting endpoint coefficient.
This proves both its convergence and
\eqref{ads:cmc:leading-integral}.
For lower degrees the integrated derivative is zero, proving their
claimed vanishing.

This order of operations uses only smooth jets on finite radial
intervals and coefficientwise endpoint limits.
In particular, convergence of the derivatives of the asymptotic
graph coefficients or of their lapse is not needed to obtain the
integral's convergence.
If the integral is zero, each of its continuous nonnegative
sphere-integrated summands is zero for every \(s>0\).
The positivity just established forces \(k=U=F=0\) there.
Smoothness gives their boundary values.
\end{proof}
